% Target: rmp-app-toric-dual
% Formula: physical spins = differences of adjacent plaquette colours
% Ink: Kernel plane frame; the blocked sheet is the house enclosure, the four
%   plaquette-colour legs leave the retained mid-side stations, and the two
%   interior diagonals recall the original lattice.
% Boundary: The four retained mid-side sites expose the plaquette-colour legs.
% Source: paper TN-Review-main.tex line 2462; paper-context-only; no author drawing
%   source
% Capabilities: kernel, plane-frame, regions, typed-ports
\tenkzkernel
\[
  % The dual blocked tensor: the four mid-side stations expose the plaquette
  % colours h_1..h_4, and the corner-to-corner diagonals are the original
  % lattice's edges crossing under the dual sheet.
  \begin{tenkz}[lattice={3x3}, frame=plane, bonds=none, boundary=none]
    \tn[at=(1,1), name=cnw, skin=none]{}
    \tn[at=(1,3), name=cne, skin=none]{}
    \tn[at=(3,1), name=csw, skin=none]{}
    \tn[at=(3,3), name=cse, skin=none]{}
    \tn[at=(1,2), name=n, skin=none]{}
    \tn[at=(2,1), name=w, skin=none]{}
    \tn[at=(2,3), name=e, skin=none]{}
    \tn[at=(3,2), name=s, skin=none]{}
    \tnwire{cnw}{cse} \tnwire{cne}{csw}
    \tnwire{n.90}{open n} \tnwire{w.180}{open w}
    \tnwire{e.0}{open e} \tnwire{s.270}{open s}
    \tnmark[form=enclosure]{(1,1) .. (3,3)}{}
  \end{tenkz}
\]
